\chapter{The First Condition for Investing}

Many people say, with complete sincerity:

\begin{quote}
When I have money, I will start investing.
\end{quote}

The sentence sounds prudent. It sounds as if the speaker respects investing enough not to rush into it. It also contains several traps.

The first trap is obvious: when, exactly, will you have money? If the answer is vague, the plan is not a plan. It is a postponement dressed as responsibility. Many people wait for the day when they will finally have enough money to begin, and that day never arrives. Their lives pass without a real beginning.

The second trap is deeper:

\begin{quote}
Investing does not require waiting until you have a large amount of money.
\end{quote}

Even that sentence is still not strong enough.

The real point is:

\begin{quote}
Investing does not necessarily require money at all in the beginning.
\end{quote}

This sounds strange only because most people confuse investing with the act of placing money into a market. That act is only one visible part of investing. Before money can be invested well, the investor must be trained. His concepts must be clear. His values must be ordered. His patience must be built. His attention must be disciplined. His mind must learn how to face uncertainty without panic or fantasy.

Money without that preparation is not capital. It is merely money exposed to danger.

\section*{Money Is Not Yet Capital}

Earlier we made an important distinction: funds become capital only when they stand behind a capable person. Money held by an unprepared mind is not yet capital, because the person behind it may lack the judgment required to protect it, allocate it, and let it grow.

Three elements turn funds into capital:

\begin{enumerate}
\item amount;
\item time horizon;
\item wisdom.
\end{enumerate}

Their order of importance is usually the reverse of what beginners imagine:

\[
\text{wisdom} > \text{time horizon} > \text{amount}
\]

The amount matters. No serious person denies that. Larger capital gives more choices, more resilience, and more room for error. But amount is not the most fundamental factor.

The core measure of investment success is not the absolute amount of money earned. It is the relative growth of the original money. A person who turns 1 into 2 has doubled his capital. A person who turns 1,000,000 into 1,100,000 has earned more in absolute terms, but the growth rate is only 10\%.

This distinction is not a mathematical decoration. It changes the mind.

\begin{center}
\begin{adjustbox}{max width=\linewidth}
\begin{tabular}{p{0.34\linewidth}p{0.46\linewidth}}
\hline
Weak measurement & Strong measurement \\
\hline
How much money did I make? & By what percentage did my capital grow? \\
Who has more money than me? & Am I improving my judgment and return quality? \\
Can I start only when the amount is large? & Can I train the investor before the amount is large? \\
\hline
\end{tabular}
\end{adjustbox}
\end{center}

If growth rate is the more essential measure, then a small amount can still teach the central lessons. If wisdom is more important than amount, then training wisdom can begin before meaningful funds appear. If time horizon is more important than amount, then the person who begins earlier has already gained something that cannot be bought later at any price.

\section*{The Fairness of Time}

Amount is not distributed fairly. Some people are born into families with capital. Some begin with debt. Some receive support, networks, houses, introductions, and safety cushions. Others begin with none of these.

This is reality. Complaining about it does not improve it.

Time is different. Not in the sentimental sense that everyone uses it equally, but in the practical sense that anyone can decide to assign a small sum, or even a simulated sum, to a long horizon. A person may not have much money, but he can decide that a small portion of his future attention and a small portion of his savings will be sentenced to a long term.

That phrase matters: sentenced to a long term.

The beginner who can take a few thousand yuan, or a few thousand dollars, or even a virtual unit in a spreadsheet, and decide not to disturb it for years has already entered the same time dimension as more advanced investors. He may not have their amount. He may not yet have their skill. But in the dimension of time horizon, he can begin standing in the right place.

\begin{center}
\begin{adjustbox}{max width=\linewidth}
\begin{tikzpicture}[font=\scriptsize, >=stealth, node distance=.9cm]
\node[draw, rounded corners=2pt, align=center, minimum width=2.2cm] (funds) {Funds};
\node[draw, rounded corners=2pt, align=center, right=of funds, minimum width=2.4cm] (amount) {Amount};
\node[draw, rounded corners=2pt, align=center, below=of amount, minimum width=2.4cm] (time) {Time\\horizon};
\node[draw, rounded corners=2pt, align=center, above=of amount, minimum width=2.4cm] (wisdom) {Wisdom};
\node[draw, rounded corners=2pt, align=center, right=1.25cm of amount, minimum width=2.4cm] (capital) {Capital};
\draw[->] (funds) -- (amount);
\draw[->] (amount) -- (capital);
\draw[->] (time) -- (capital);
\draw[->] (wisdom) -- (capital);
\end{tikzpicture}
\end{adjustbox}
\end{center}

Wisdom is the most interesting element because it feels difficult and simple at the same time. It is difficult because no one can give you mature judgment instantly. It is simple because no one has a permanent birthright over it. Investment wisdom is not inherited in the way a house can be inherited. It must be learned, practiced, corrected, and deepened.

In that sense, no one has an unbeatable innate advantage.

This should be both comforting and frightening. Comforting, because you are not disqualified by birth. Frightening, because if you do not train, the responsibility remains yours.

\section*{Do Not Compete With Other Investors}

Once people hear that growth rate matters, some immediately want to compete. They want to compare returns, portfolios, predictions, and cleverness.

This impulse is mostly useless.

Investing is your own affair. Each person begins with a different amount, a different family situation, a different income, a different risk tolerance, a different knowledge base, and a different time horizon. Comparison introduces emotional noise where clear judgment is needed. It makes the small investor arrogant after a lucky gain and ashamed after a temporary loss. It makes him imitate people whose circumstances he does not understand.

The investor you must beat is yesterday's self.

This is not a slogan of comfort. It is an operational rule. Your goal is to think more clearly than before, record more consistently than before, understand risk better than before, and behave with more stability than before. If you improve those things, the visible results will eventually have a better foundation. If you do not improve those things, even a lucky result may make you more dangerous to yourself.

The market does not need your self-esteem. It does not care whether you feel behind. It rewards and punishes behavior. Therefore, attention should not be spent on comparison unless comparison produces a concrete method.

\section*{The Question Behind Every Decision}

The most useful thinking method is almost embarrassingly plain:

\begin{quote}
What is the most important thing here?
\end{quote}

Ask it again and again.

What is most important in investing? Not the excitement of the market. Not the gossip around a company. Not the amount someone else already has. The most important thing is the investor's wisdom: the ability to understand, decide, wait, review, and correct.

What is most important in defining capital? Amount, time horizon, and wisdom. Among these, wisdom is first.

What is most important in measuring investment results? Relative growth, not merely absolute money.

What is most important when you feel you cannot start? Finding a form of practice that begins immediately.

This question is the working tool of values. A person without values drifts. A person with vague values still drifts, only with better language. A person who repeatedly asks what is more important begins to sharpen the hierarchy inside his mind. That hierarchy then determines action.

Many people fail not because they lack desire, but because their concepts are unclear and their values are disordered. They want wealth freedom, but they cannot define wealth. They want investment returns, but they do not know what counts as success. They say they will start when they have money, but they cannot say what amount qualifies as ``having money.''

If the concept is foggy, the decision will be foggy.

\section*{Wealth Freedom Is Not the Final Goal}

There is another value error hidden in the phrase ``I want wealth freedom.''

Wealth freedom is important, but it is not the ultimate purpose of a life. It is a milestone. For some people, it may not even feel like a milestone once reached; it may simply be the byproduct of doing valuable things in the right way for a long time.

Look at people who actually build serious wealth. Many of them are not aiming only at escape. They are building companies, solving problems, improving systems, creating products, teaching, researching, organizing, or contributing. They have something more important than money in front of them, and that ``more important'' thing is often what makes money possible.

At minimum, one thing should be placed above the desire for money:

\begin{quote}
the continuous refinement of one's concepts and values.
\end{quote}

If the mind is full of leaks, more money only increases the speed of collapse. Every thinking leak eventually breaks through the dam. The more capital behind the leak, the more spectacular the damage.

This is why the training must begin before the money arrives. If a person waits until he has substantial funds before he begins correcting his thinking, he has waited too long. The correction should begin when the stakes are small, when the cost of error is still educational rather than destructive.

\section*{A One-Unit Investment Practice}

So what should a person do if he truly has no money to invest?

Create a simple spreadsheet.

Choose one company that you believe, with reasons, may continue to grow over a long period. It does not have to be one of the famous technology companies discussed earlier. The company itself is not sacred. If you can read English comfortably, you may choose an overseas public company. If not, choose a domestic company whose reports and news you can actually follow.

Then imagine that you bought exactly one unit of it: one dollar, one yuan, or the smallest convenient unit in your currency.

At the end of each month, record the price and calculate the percentage change from your original one-unit purchase.

That is all.

\begin{center}
\begin{adjustbox}{max width=\linewidth}
\begin{tabular}{p{0.22\linewidth}p{0.22\linewidth}p{0.22\linewidth}p{0.20\linewidth}}
\hline
Date & Company & Price & Change \\
\hline
Start & Example Co. & 1.00 & 0\% \\
Month 1 & Example Co. & 1.03 & 3\% \\
Month 2 & Example Co. & 0.97 & -3\% \\
Month 3 & Example Co. & 1.08 & 8\% \\
\hline
\end{tabular}
\end{adjustbox}
\end{center}

The amount must be one unit because the exercise is not about pretending to be rich. It is about training the habit of thinking in relative terms. One yuan, one dollar, one thousand, and ten million all share the same growth-rate logic. From the beginning, the mind must learn to ask:

\begin{quote}
What happened to the capital in percentage terms?
\end{quote}

You may track several companies, but do not track too many. Three to five is enough for a beginner. More than that may look ambitious, but it often becomes scattered attention. Even simulated investing requires real mental work.

The rule is strict: update the sheet only once a month.

Do not check the price every day. Do not check it every hour. Do not ``just take a look.'' Looking too often is not harmless. It trains the wrong reflex. It turns investing into stimulation. It builds the habit of emotional reaction to short-term movement.

If you break the rule and check early, impose a small penalty on yourself. The point is not punishment for its own sake. The point is to teach the mind that attention has rules and that rules matter.

\section*{Why the Exercise Is Real}

Many people will ask, ``What is the point of this?''

The honest answer is: if you could already see the point deeply, you would probably have started some form of disciplined practice long ago.

The exercise looks small because beginners underestimate the training it contains.

First, it trains patience. Updating a number once a month seems easy, but it is difficult for people whose sense of time is measured in hours and days. Investing requires the mind to perceive weeks, months, years, and decades. Without that expanded time sense, volatility feels larger than it is, and patience remains a word rather than a capacity.

Second, it trains attention. During the month, you will be tempted to check. Resisting that temptation is part of the practice. You are learning that not every available piece of data deserves your attention. You are learning that attention should serve a chosen system.

Third, it trains research. If you track a company, you will naturally begin to read its news, products, reports, competitors, and industry. The company becomes a living case study. You are no longer reading abstract business language. You are watching a real organization meet reality.

Fourth, it may train English. If you follow overseas companies, you will have a reason to read English reports and announcements. Learning a language by ``studying'' it often becomes lifeless. Using it to understand something you care about gives the mind a reason to stay awake. Use the language, and the language gradually becomes less frightening.

Fifth, it trains record keeping. A mind that cannot record cannot review. A person who cannot review cannot improve reliably. Memory is too self-serving. Records are less polite and more useful.

\section*{Twelve Months Before Entry}

Continue the practice for at least twelve months.

Only then can you call yourself a beginner.

This may sound slow. In fact, it is fast. Many people enter markets with no training at all. They believe they have entered an investment field, but in practice they have opened the door to a casino. They bring impatience, greed, fear, pride, hearsay, and untested assumptions, then wonder why the market feels cruel.

The point of a year of practice is not to predict prices correctly. The point is to reveal your own mind.

You will see whether you can follow a rule. You will see whether you become excited after a rise and depressed after a fall. You will see whether you secretly want to change the rules to protect your ego. You will see whether you can read calmly when the price disagrees with your expectation. You will see whether you can keep a small commitment when no one is watching.

This is investing education at the correct scale. Small enough not to ruin you. Real enough to expose you.

\section*{Becoming Worthy of the Opportunity}

Charlie Munger's line is worth keeping close:

\begin{quote}
The best way to get something is to deserve it.
\end{quote}

Most people focus on what they want. Fewer ask what they have to exchange for it. Life is full of exchange. If you do not yet receive what you want, one useful question is not ``Why is the world unfair?'' but ``What must I improve so that I can make a better exchange with the world?''

In investing, the question becomes:

\begin{quote}
What kind of person is worthy of this opportunity?
\end{quote}

A worthy investor is not someone who never makes mistakes. That person does not exist. A worthy investor is someone who can define terms clearly, respect time, measure growth correctly, control attention, record decisions, review errors, avoid unnecessary comparison, and keep improving his judgment.

The whole practice of this book is to close mental leaks before capital grows large. When the real opportunity appears, the unprepared person asks for instructions, hesitates, reacts emotionally, or enters blindly. The prepared person may still be uncertain, but he has a system.

\begin{center}
\begin{adjustbox}{max width=\linewidth}
\begin{tikzpicture}[font=\scriptsize, >=stealth, node distance=.8cm]
\node[draw, rounded corners=2pt, align=center, minimum width=2.2cm] (concepts) {Clear\\concepts};
\node[draw, rounded corners=2pt, align=center, right=of concepts, minimum width=2.2cm] (values) {Ordered\\values};
\node[draw, rounded corners=2pt, align=center, right=of values, minimum width=2.2cm] (practice) {Monthly\\practice};
\node[draw, rounded corners=2pt, align=center, right=of practice, minimum width=2.2cm] (judgment) {Better\\judgment};
\node[draw, rounded corners=2pt, align=center, right=of judgment, minimum width=2.2cm] (capital) {Capital\\ready};
\draw[->] (concepts) -- (values);
\draw[->] (values) -- (practice);
\draw[->] (practice) -- (judgment);
\draw[->] (judgment) -- (capital);
\end{tikzpicture}
\end{adjustbox}
\end{center}

The first condition for investing is therefore not money.

It is willingness.

Willingness to begin before the beginning looks impressive. Willingness to train when the amount is small. Willingness to ask what is most important. Willingness to think in ratios. Willingness to record. Willingness to wait. Willingness to discover that what you thought you understood was only the first surface of understanding.

A year from now, you can have twelve months of experience or twelve more months of excuses. The required tool is simple. The required money can be imaginary. The required wisdom begins with one decision:

\begin{quote}
Start the practice now.
\end{quote}
